\documentclass[11pt]{article}
\usepackage[a4paper,margin=1.9cm]{geometry}
\usepackage{fontspec}
\setmainfont{Latin Modern Roman}
\usepackage{titlesec}
\usepackage{enumitem}
\usepackage[table]{xcolor}
\usepackage{longtable}
\usepackage{booktabs}
\usepackage{array}
\usepackage{pifont}
\usepackage{amssymb}
\usepackage{tikz}
\usepackage{tcolorbox}
\usepackage{hyperref}
\usepackage{fancyhdr}
\usepackage{parskip}

\definecolor{gxteal}{HTML}{0F766E}
\definecolor{gxcyan}{HTML}{0284C7}
\definecolor{gxslate}{HTML}{334155}
\definecolor{gxamber}{HTML}{B45309}
\definecolor{gxred}{HTML}{B91C1C}
\definecolor{gxgreen}{HTML}{15803D}
\definecolor{gxviolet}{HTML}{6D28D9}

\hypersetup{colorlinks=true, linkcolor=gxteal, urlcolor=gxcyan,
  pdftitle={GlobeX -- PS 26132 and PS 26102 Honest Build Assessment}, pdfauthor={GlobeX Team}}

\titleformat{\section}{\large\bfseries\color{gxteal}}{\thesection}{0.6em}{}
\titleformat{\subsection}{\normalsize\bfseries\color{gxslate}}{\thesubsection}{0.6em}{}
\titlespacing*{\section}{0pt}{1.7em}{0.7em}

\pagestyle{fancy}
\fancyhf{}
\renewcommand{\headrulewidth}{0.4pt}
\fancyhead[L]{\small\color{gxslate}GlobeX --- PS 26132 \& PS 26102 Honest Assessment}
\fancyhead[R]{\small\color{gxslate}SIH 2026}
\fancyfoot[C]{\small\thepage}

\newcommand{\yes}{{\color{gxgreen}\ding{108}}}
\newcommand{\pmid}{{\color{gxcyan}\ding{109}}}
\newcommand{\non}{{\color{black!22}\ding{109}}}
\newcommand{\bld}{{\color{gxamber}\ding{53}}}

\newtcolorbox{callout}[1]{colback=gxteal!5, colframe=gxteal!55, boxrule=0.6pt, arc=2pt,
  left=6pt, right=6pt, top=5pt, bottom=5pt, fonttitle=\bfseries\small, coltitle=white,
  colbacktitle=gxteal, title={#1}}
\newtcolorbox{uspbox}[1]{colback=gxviolet!5, colframe=gxviolet!55, boxrule=0.6pt, arc=2pt,
  left=6pt, right=6pt, top=5pt, bottom=5pt, fonttitle=\bfseries\small, coltitle=white,
  colbacktitle=gxviolet, title={#1}}
\newtcolorbox{warnbox}[1]{colback=gxamber!6, colframe=gxamber!65, boxrule=0.6pt, arc=2pt,
  left=6pt, right=6pt, top=5pt, bottom=5pt, fonttitle=\bfseries\small, coltitle=white,
  colbacktitle=gxamber, title={#1}}
\newtcolorbox{redbox}[1]{colback=gxred!6, colframe=gxred!65, boxrule=0.6pt, arc=2pt,
  left=6pt, right=6pt, top=5pt, bottom=5pt, fonttitle=\bfseries\small, coltitle=white,
  colbacktitle=gxred, title={#1}}

\begin{document}

% ================= TITLE =================
\begin{titlepage}
\centering
\vspace*{1.6cm}
{\Huge\bfseries\color{gxteal} PS 26132 \& PS 26102}\\[0.4em]
{\Huge\bfseries\color{gxteal} The Honest Build Assessment}\\[1.8em]
{\Large\color{gxslate} What's genuinely strong, what's genuinely weak,}\\[0.3em]
{\Large\color{gxslate} what could sink the demo, and how much is realistically buildable}\\[2.6em]
{\normalsize\color{gxslate} 26132 --- Market linkages \& price discovery for farmers (Govt.\ of Maharashtra / MSInS)}\\[0.2em]
{\normalsize\color{gxslate} 26102 --- Anomaly, fraud \& inefficiency detection in MPLADS (MoSPI)}\\[0.5em]
{\normalsize\color{gxslate} Prepared for GlobeX --- 2 September 2026}\\[2.2em]
\begin{tabular}{cc}
\textbf{\color{gxteal}\Large 26132} & \textbf{\color{gxteal}\Large 26102} \\
{\small fit rank 1 / 229 \ \ coverage 9/9} & {\small fit rank 31 / 229 \ \ coverage 9/9} \\[0.3em]
{\small 4 gaps to close} & {\small 3 gaps to close} \\
\end{tabular}
\vfill
{\footnotesize\color{gxslate} Internal working document --- not for external distribution as-is}
\end{titlepage}

% ================= VERDICT =================
\section*{The one-line verdict on each}
\begin{callout}{26132 --- build it, but close the gaps first}
The best-measured fit GlobeX has against any of the 229 statements, and the escrow layer answers the exact
fear (``payment reliability'') the problem statement names in its own text. Three specific things are not
demo-ready yet: a stub data path inside counterparty matching, a forecaster whose headline number
contradicts its own benchmark file, and an escrow flag that is off by default.
\end{callout}
\begin{callout}{26102 --- the strongest match to our best-validated asset, with one hard rule}
Ranks 31st on raw similarity but is the closest match in the whole dataset to GlobeX's single strongest,
most rigorously validated component: the anomaly/risk stack. The entire pitch stands or falls on one
discipline: never present the 0.98 F1 trained on trade data as a result on MPLADS data. Retrain on the
scheme's own public data and report that number, whatever it is.
\end{callout}

% ================= GOOD =================
\clearpage
\section{What's genuinely good}

\subsection*{PS 26132 --- Market linkages \& price discovery}
Across a four-profile TF-IDF similarity sweep run against all 229 statements, 26132 ranks first on the
profile that matters most --- product shape --- by a wide margin: 65\% clear of whatever sits second. Every
one of the nine things the problem statement asks for (price visibility, sale-window advice, verified
matching, quality/grading, lot creation, logistics, payment tracking, dispute handling, transparent
records) maps to something GlobeX already has a working implementation of.

Most competing teams will build a price dashboard with a buyer directory. The problem statement itself
names ``payment reliability'' as the thing farmers actually fear, and a dashboard does not fix that ---
protecting the transaction does. The escrow contract already implements fund-on-acceptance, condition-gated
release, a dispute state that freezes funds instead of losing them, and a write-intent-before-chain pattern
that never marks a transfer complete without a confirmed on-chain receipt --- the kind of engineering
discipline a technical judge notices.

The MCDA destination-ranking engine (explainable, 8 weighted dimensions, backtested at $\rho = 0.884$,
precision@5 of 0.8 on a real 48{,}445-row dataset) becomes the sale-window recommender almost without
modification. And the one dataset this idea absolutely needs --- mandi price and arrival data --- is
genuinely public and reachable today through a documented API (data.gov.in / AIKosh), not locked behind a
request form.

\subsection*{PS 26102 --- Anomaly detection in MPLADS}
This is the closest match in the entire dataset to GlobeX's strongest asset. The trade anomaly classifier
is trained, chronologically validated, and performing at 0.98 F1; the unsupervised risk ensemble (Isolation
Forest + GRU autoencoder) exists precisely because trade data has no labelled fraud --- exactly the
condition public-fund irregularity data is in. Sanctioned-work expenditure records are structurally the
same object as trade transactions (an amount, a date, a category, a counterparty), so the rolling-statistics
and peer-deviation features transfer with minimal rework, and TreeSHAP already attaches an explanation to
every flag --- non-negotiable when the subject is public money.

The USP here is unusually strong and hard to copy: anchoring the flag itself (model version, input snapshot
hash, score, officer disposition) to the evidence ledger, so the oversight trail is as tamper-evident as the
spending it oversees. Paired with a deliberate refusal to equate ``anomalous'' with ``fraudulent'' --- the
existing risk documentation already makes this distinction --- the pitch is restrained in exactly the way a
domain involving public accusation needs it to be.

MPLADS data is also genuinely public: the scheme runs its own dashboard (mplads.gov.in, and a newer MoSPI
dashboard at mplads.mospi.gov.in --- run by the same division that owns this PS), data.gov.in carries
MPLADS-tagged datasets, and third-party aggregators like Dataful publish cleaned state/constituency-wise
fund-release and expenditure figures. Getting real numbers to train and demo against is not the blocker
here.

% ================= BAD =================
\clearpage
\section{What's genuinely weak}

\subsection*{PS 26132}
\renewcommand{\arraystretch}{1.3}
\begin{longtable}{@{}p{4.1cm}p{7.6cm}p{3.6cm}@{}}
\toprule
\textbf{Weakness} & \textbf{What's actually true} & \textbf{Why it matters} \\
\midrule
\endhead
Counterparty matching has a stub path & \texttt{counterparty\_controller.py} still contains an
\texttt{hashlib.md5}-seeded fake-data generator on at least part of its matching path. A separate, genuine
endpoint (company directory: TF-IDF + cosine + valuation blending) is real but not what ``verified buyers''
would naturally route through. & A judge who asks for a live match on real-looking data can hit the fake
path by accident. \\[0.4em]
No public dataset for FPOs or buyer credentials & No equivalent of the Agmarknet price feed exists for
farmer-producer-organisation directories or buyer verification records. & The matching engine's ``ground
truth'' has to be synthesised, and unlabelled synthetic data looks exactly like a stub-data accusation from
the outside. \\[0.4em]
Forecaster headline number contradicts its own benchmark & The README claims 8.42\% MAPE; the repo's own
\texttt{benchmark\_comparison.csv} shows the deployed Dual-Head GRU actually scoring 54.49\% MAPE on demand
--- worse than a plain moving average (29.42\%). & The single most dangerous line in the project if a judge
cross-reads the README against the repo. \\[0.4em]
Escrow is off by default & \texttt{ESCROW\_ENABLED} defaults to \texttt{false}, and only one migration file
exists in the whole repository --- the schema is not version-controlled. & The best differentiator in the
project is a flag flip and an untested schema away from actually running live. \\
\bottomrule
\end{longtable}

\subsection*{PS 26102}
\renewcommand{\arraystretch}{1.3}
\begin{longtable}{@{}p{4.1cm}p{7.6cm}p{3.6cm}@{}}
\toprule
\textbf{Weakness} & \textbf{What's actually true} & \textbf{Why it matters} \\
\midrule
\endhead
The headline accuracy is from the wrong domain & 0.98 F1 was measured on trade-anomaly data, never on
MPLADS expenditure data. There is currently no retrained result on the actual scheme data. & This is the
single biggest risk on this PS. Presenting 0.98 as if it were an MPLADS result is a conflation a judge will
catch immediately, and it will discount everything else in the pitch. \\[0.4em]
Work-level detail is not cleanly structured & The public dashboards give fund-release and expenditure
totals well; individual sanctioned-work records with cost estimates and execution status are less uniformly
structured and partly locked in per-constituency HTML tables rather than a clean bulk API. & Several of the
PS's specific asks (cost-estimate deviation, duplicate/overlapping works) need work-level granularity, which
takes real scraping/cleaning effort to get to. \\[0.4em]
No ground truth for fraud, by design & There is no labelled ``this was fraud'' dataset for public works,
just as there is none for trade --- which is why the unsupervised ensemble exists, but it also means there
is no way to report a precision/recall number against real fraud, only against synthetic or proxy labels. &
Anomaly $\neq$ fraud has to be said out loud in the pitch, not discovered by the judges; otherwise it reads
as an accuracy gap rather than a deliberate, defensible design choice. \\
\bottomrule
\end{longtable}

\begin{redbox}{The two things to fix before anyone else sees either pitch}
For 26132: correct the README's forecasting claim now --- either retrain and report honestly, or switch the
headline number to the XGBoost quantile forecaster's real performance and drop the GRU MAPE figure entirely.
For 26102: retrain the anomaly/risk stack on real MPLADS data (even a modest slice) before finals, and
report that number as the headline, not the 0.98 trade figure. Both are a handful of focused hours, not new
capability.
\end{redbox}

% ================= CHALLENGES =================
\clearpage
\section{Challenges, and exactly how to tackle each one}

\subsection*{PS 26132}
\renewcommand{\arraystretch}{1.35}
\begin{longtable}{@{}p{3.7cm}p{4.1cm}p{7.0cm}@{}}
\toprule
\textbf{Challenge} & \textbf{Risk if left alone} & \textbf{How to tackle it} \\
\midrule
\endhead
Government price data is messy --- inconsistent commodity names, missing days, non-standard units & Dirty
ingestion makes the forecaster look wrong live & Build a normalisation layer early: canonical
commodity/unit mapping, a gap-fill policy, a validation pass that flags suspicious jumps before they reach
the forecaster. \\[0.4em]
Stub logic inside counterparty matching & Live demo can hit fake data by accident & Route ``verified
buyer'' explicitly through the real company-directory endpoint; if the stub path stays for anything, label
it unmistakably as a placeholder. \\[0.4em]
No real FPO/buyer directory & Matching demo runs on synthetic data that looks bad if presented as real &
Build a clearly-labelled pilot seed set (50--150 records styled on real APMC/cooperation structures) and
say so in the pitch: ``pilot seed data; onboarding real FPOs is the first 90-day milestone.'' \\[0.4em]
Forecaster overclaiming / domain shift & GRU was validated on trade-corridor data, not mandi series &
Lead with the XGBoost quantile forecaster (F2), recalibrated on actual ingested mandi series; present every
recommendation as a range with a confidence band, never a bare number. \\[0.4em]
Escrow reliability under demo conditions & A live chain call failing in front of judges is worse than never
having built it & Flip \texttt{ESCROW\_ENABLED} on, dry-run the happy path and dispute path repeatedly on
testnet, add schema migrations now, keep a recorded fallback run. \\
\bottomrule
\end{longtable}

\subsection*{PS 26102}
\renewcommand{\arraystretch}{1.35}
\begin{longtable}{@{}p{3.7cm}p{4.1cm}p{7.0cm}@{}}
\toprule
\textbf{Challenge} & \textbf{Risk if left alone} & \textbf{How to tackle it} \\
\midrule
\endhead
No MPLADS-trained result yet & The only number in hand (0.98 F1) is from the wrong domain & Pull a real
slice of MPLADS expenditure/works data (mplads.gov.in, MoSPI dashboard, data.gov.in, Dataful), engineer the
same rolling-stats/peer-deviation features, retrain, and report whatever number results --- even if it is
far lower than 0.98. A modest honest number here beats a borrowed impressive one. \\[0.4em]
Work-level records are semi-structured & Cost-estimate deviation and duplicate-work detection need
per-work granularity that isn't cleanly exposed in bulk & Scope the demo to a small number of constituencies
where per-work detail can realistically be scraped and cleaned in the time available, rather than
attempting national coverage. \\[0.4em]
Anomaly vs.\ fraud framing & An unprepared team either overclaims (``we detect fraud'') or looks weak when
pressed on accuracy against real fraud & State the distinction proactively in the pitch: the system flags
statistical irregularity for human review; it does not accuse. This is a documented design position, not an
improvised excuse. \\[0.4em]
Sensitivity of accusing public officials/agencies & A false or careless flag against a real MP's fund or
agency is reputationally serious, unlike a commercial trade anomaly & Keep every flag routed to a
human-in-the-loop disposition step, never auto-publish a flag, and make the evidence-ledger anchoring cover
the review decision too, not just the raw score. \\
\bottomrule
\end{longtable}

% ================= SCOPE =================
\clearpage
\section{How much can realistically be built --- a scope plan}

\subsection*{PS 26132}
\textbf{Must-have:} ingestion of Agmarknet/AIKosh mandi data for 2--3 states and a handful of commodities,
cleaned and normalised; sale-window recommendation via the XGBoost quantile forecaster presented as a price
band; buyer matching through the real (non-stub) endpoint against the pilot seed dataset, with a visible
trust score; escrow happy path live on testnet with the evidence-ledger entry visible; a clickable front-end
flow start to finish.

\textbf{Should-have:} dispute path demoed live, not just described; document/grading verification against a
sample certificate and lot record; the audit trail exposed in the UI; coverage widened to 4--5 states once
the core flow is solid.

\textbf{Could-have (genuine stretch, do not commit):} an image-based grade classifier if a labelled photo
dataset can be assembled in time; a vernacular or WhatsApp-bot front end; multi-buyer competitive offers.

\textbf{Won't-have --- say so upfront:} live eNAM transaction data or GSTN/bank payment-rail integration;
nationwide coverage at launch; a trained (rather than rules-based) grading model unless the stretch goal is
actually reached.

\subsection*{PS 26102}
\textbf{Must-have:} a real, cleaned MPLADS dataset for a handful of constituencies (fund release,
expenditure, work-level records where scrapeable); the anomaly classifier and unsupervised risk ensemble
retrained on that data with an honestly reported score; TreeSHAP-style explanation attached to every flag; a
monitoring dashboard showing flagged works with their contributing factors; the flag-anchoring-to-ledger
mechanism working end to end for at least one flagged case.

\textbf{Should-have:} duplicate/overlapping-work detection via TF-IDF similarity across constituencies;
cost-estimate peer-deviation scoring; a human disposition workflow (reviewer marks a flag reviewed/dismissed,
and that decision is itself anchored).

\textbf{Could-have (genuine stretch):} coverage extended across more states/constituencies; a
trend/early-warning view that aggregates flags by implementing agency over time.

\textbf{Won't-have --- say so upfront:} any claim of detecting or proving fraud (only irregularity worth
review); real-time ingestion tied to live government systems rather than periodically refreshed public data;
national-scale coverage at launch.

\begin{warnbox}{On ``how much can we make'' as a competitive question, for both}
Conditional on closing the gaps listed above, both are top-tier entries on paper: 26132 for the measured
\#1 fit and a differentiator the problem statement names outright; 26102 for pairing GlobeX's single
best-validated component with an accountability mechanism (ledger-anchored flags) that is genuinely hard to
copy in a hackathon window. Left unfixed, in both cases the failure mode is identical --- a sharp judge finds
the domain-mismatched number or the stub path in the first ten minutes of questions, and a top-tier
submission reads as an average one. The idea and the fit do not change; the trust does. In both cases the
gap between those two outcomes is single-digit hours of focused, unglamorous work, not new capability.
\end{warnbox}

\section*{Bottom line}
26132 is the stronger overall submission --- best fit in the dataset, and a differentiator the PS names in
its own text. 26102 is the stronger showcase for GlobeX's actual ML rigor and pairs unusually well with the
blockchain layer through ledger-anchored flags, but it carries one hard, non-negotiable discipline: never
show the trade-trained 0.98 F1 as an MPLADS result. Whichever is filed, walk in having already fixed the
domain-mismatched numbers, rerouted matching or retrained on real data as applicable, and turned on and
tested the chain layer --- before touching a single new feature.

\end{document}
